\section{Προτασιακή Λογική}
\paragraph{Βασικές έννοιες για τις προτάσεις}
Μια πρόταση στα μαθηματικά συμβολίζεται συνήθως με κεφαλαίο γράμμα π.χ. $A,B,\ldots$. Για παράδειγμα η 
\begin{center}
$A$: \textit{Ο αριθμός $3$ είναι ακέραιος}
\end{center}
αποτελεί μια μαθηματική πρόταση. Ένα θεώρημα στα μαθηματικά αποτελείται από δύο προτάσεις, την \textbf{υπόθεση} $A$ και το \textbf{συμπέρασμα} $B$, που συνδέονται με την πράξη της συνεπαγωγής. Έχει δηλαδή τη μορφή
\[ A\Rightarrow B \]
και έχει το νόημα ότι \bmath{αν ισχύει η $A$ τότε ισχύει η $B$}. Για τις προτάσεις ισχύουν τα εξής:
\begin{itemize}
\item Κάθε πρόταση μπορεί να θεωρηθεί ως μεταβλητή και η μεταβλητή αυτή παίρνει δύο τιμές: \textbf{Αληθής} (\textcolor{check}{\faCheck}) και \textbf{Ψευδής} (\textcolor{error}{\faTimes}).
\item Η πράξη της \textbf{άρνησης} στις προτάσεις, μας δίνει το αντίθετο μιας πρότασης, συμβολίζεται με $\neg$ και μετατρέπει μια πρόταση από αληθή σε ψευδή και αντίστροφα. Για παράδειγμα 
\begin{center}
$\neg A$: \textit{Ο αριθμός $3$ \textbf{δεν} είναι ακέραιος}.
\end{center}
\item Το \textbf{αντίστροφο} μιας συνεπαγωγής $A\Rightarrow B$ γράφεται είτε $A\Leftarrow B$ είτε $B\Rightarrow A$.
\item Για μια συνεπαγωγή $A\Rightarrow B$, ο συνδυασμός της άρνησης και του αντιστρόφου ονομάζεται \textbf{αντιθετοαντίστροφη} συνεπαγωγή και συμβολίζεται $\neg B\Rightarrow\neg A$. Έχει την έννοια του \bmath{αν δεν ισχύει το συμπέρασμα $B$ τότε δεν ισχύει και η υπόθεση $A$}.
\item Η \textbf{διπλή συνεπαγωγή}, ή αλλιώς \textbf{ισοδυναμία}, συμβολίζεται $\Leftrightarrow$ και δηλώνει ότι ισχύει και το ορθό $A\Rightarrow B$ και το αντίστροφο $A\Leftarrow B$. Όταν δύο προτάσεις $A,B$ είναι ισοδύναμες τότε είναι είτε και οι δύο αληθείς είτε και οι δύο ψευδείς.
\end{itemize}
\paragraph{Θεωρήματα Γ' Λυκείου}
Θα χρησιμοποιήσουμε τώρα όσα αναφέραμε στα θεωρήματα και τι προτάσεις των μαθηματικών προσανατολισμού της Γ' Λυκείου. Ο γενικός σκοπός μας είναι να διατυπώσουμε για κάθε θεώρημα,  τις εξής $4$ «μορφές» που μπορεί να πάρει και να τις μελετήσουμε ως προς την ισχύ:
\begin{center}
\begin{mytblr}{}
\circled{1}: Θεώρημα & \SetCell{orange!90!black}\circled{2}: Αντίστροφο & \circled{3}: Αντιθετοαντίστροφο του \circled{1} & \SetCell{orange!90!black}\circled{4}: Αντιθετοαντίστροφο του \circled{2}\\
$A\Rightarrow B$ & \SetCell{orange!50!white} $B\Rightarrow A$ & $\neg B\Rightarrow \neg A$ & \SetCell{orange!50!white}$\neg A\Rightarrow\neg B$
\end{mytblr}
\end{center}
Θα χρησιμοποιήσουμε επίσης τον εξής σημαντικό κανόνα:
\begin{center}
\textbf{Κάθε συνεπαγωγή είναι ισοδύναμη με την αντιθετοαντίστροφή της}
\end{center}
ο οποίος μας δίνει τα ακόλουθα συμπεράσματα:
\begin{center}
\circled{1}$\Leftrightarrow$\circled{3}\ \ \text{και}\ \ \circled{2}$\Leftrightarrow$\circled{4}
\end{center}
\newpage
\protash{Μονοτονία και \bmath{$1-1$}}

Αν μια συνάρτηση $f$ είναι γνησίως μονότονη τότε είναι και $1-1$.

\begin{center}
\begin{mytblr}{}
\circled{1}: $A\Rightarrow B$ & \SetCell{orange!90!black}\circled{2}: $B\Rightarrow A$ & \circled{3}: $\neg B\Rightarrow \neg A$ & \SetCell{orange!90!black}\circled{4}: $\neg A\Rightarrow\neg B$\\
\true & \SetCell{orange!50!white}\false & \true & \SetCell{orange!50!white}\false
\end{mytblr}
\end{center}

\theor{\true}
Αν η $f$ είναι γνησίως μονότονη $\Rightarrow$ Η $f$ είναι $1$-$1$.\\
{\small Δηλαδή: διαφορετικά $x$ δίνουν πάντα διαφορετικές τιμές.}\\\\
\antistr{\false}
Αν η $f$ είναι $1$-$1$ $\Rightarrow$ Η $f$ είναι γνησίως μονότονη.\\
\textcolor{error}{\faTimes\ \textbf{Αντιπαράδειγμα:}} Η $f(x) = \frac{1}{x}$, $x\in\mathbb{R}^*$, είναι $1$-$1$ αλλά ούτε αύξουσα ούτε φθίνουσα σε ολόκληρο το $\mathbb{R}^*$.

\aath{\true}
Αν η $f$ \textbf{δεν} είναι $1$-$1$ $\Rightarrow$ Η $f$ \textbf{δεν} είναι γνησίως μονότονη.\\
{\small Ισοδύναμο με το θεώρημα.}\\\\
\aaa{\false}
Αν η $f$ \textbf{δεν} είναι γνησίως μονότονη $\Rightarrow$ Η $f$ \textbf{δεν} είναι $1$-$1$.\\
\textcolor{error}{\faTimes\ \textbf{Αντιπαράδειγμα:}} Ίδιο: $f(x)=\frac{1}{x}$.

\vspace{1cm}

%% ===================== Θεώρημα Bolzano =====================
\protash{Θεώρημα \eng{Bolzano}}

Θεωρούμε μια συνάρτηση $f$ ορισμένη σε ένα κλειστό διάστημα $[a,\beta]$. Αν
\begin{itemize}
\item η $f$ είναι συνεχής στο κλειστό διάστημα $[a,\beta]$ και 
\item $f(a)\cdot f(\beta)<0$
\end{itemize}
τότε θα υπάρχει τουλάχιστον ένας αριθμός $ x_0\in(a,\beta) $ έτσι ώστε να ισχύει $ f(x_0)=0 $.
\begin{center}
\begin{mytblr}{}
\circled{1}: $A\Rightarrow B$ & \SetCell{orange!90!black}\circled{2}: $B\Rightarrow A$ & \circled{3}: $\neg B\Rightarrow \neg A$ & \SetCell{orange!90!black}\circled{4}: $\neg A\Rightarrow\neg B$\\
\true & \SetCell{orange!50!white}\false & \true & \SetCell{orange!50!white}\false
\end{mytblr}
\end{center}

\theor{\true}
Αν η $f$ είναι \textbf{συνεχής} στο $[a,\beta]$ και ισχύει $f(a)\cdot f(\beta)<0$ $\Rightarrow$ Η εξίσωση $f(x)=0$ έχει τουλάχιστον μία ρίζα στο $(a,\beta)$.\\\\
\antistr{\false}
{\kerkissans 1η μορφή:} Αν η $f$ είναι συνεχής και υπάρχει $x_0\in(a,\beta)$ τέτοιο ώστε $f(x_0)=0$ $\Rightarrow$ Ισχύει $f(a)\cdot f(\beta)<0$.\\
\textcolor{error}{\faTimes\ \textbf{Αντιπαράδειγμα:}} $f(x)=x^2$ στο $[-1,1]$: $f(0)=0$ αλλά $f(-1)\cdot f(1)=1>0$.\\\\
{\kerkissans 2η μορφή:} Αν υπάρχει $x_0\in(a,\beta)$ τέτοιο ώστε $f(x_0)=0$ και ισχύει $f(a)\cdot f(\beta)<0$ $\Rightarrow$ Η $f$ είναι συνεχής.\\
\textcolor{error}{\faTimes\ \textbf{Πρόβλημα:}} Η ρίζα και η αλλαγή προσήμου δεν εγγυώνται συνέχεια.\\\\
{\kerkissans 3η μορφή:} Αν ισχύει $f(a)\cdot f(\beta)<0$ χωρίς η $f$ να είναι συνεχής $\Rightarrow$ Υπάρχει τουλάχιστον μία ρίζα της εξίσωσης $f(x)=0$.\\
\textcolor{error}{\faTimes\ \textbf{Αντιπαράδειγμα:}} $f(x)=\begin{cases} -1, & x\le 0 \\ 1, & x>0 \end{cases}$ στο $[-1,1]$: $f(-1)\cdot f(1)<0$ αλλά η $f$ δεν μηδενίζεται!\\\\
\aath{\true}
Αν η $f$ είναι \textbf{συνεχής} στο $[a,\beta]$ και η εξίσωση $f(x)=0$ \textbf{δεν} έχει λύση στο $(a,\beta)$
$\Rightarrow$ Ισχύει $f(a)\cdot f(\beta)\ge 0$.\\
{\small Αν δεν μηδενίζεται και είναι και συνεχής, τότε διατηρεί σταθερό πρόσημο.}\\\\
\aaa{\false}
Αν η $f$ είναι \textbf{συνεχής} στο $[a,\beta]$ και ισχύει $f(a)\cdot f(\beta)\ge 0$ $\Rightarrow$ Η εξίσωση $f(x)=0$ \textbf{δεν} έχει λύση στο $(a,\beta)$.\\
\textcolor{error}{\faTimes\ \textbf{Αντιπαράδειγμα:}} $f(x)=(x-1)^2$ στο $[0,2]$: $f(0)\cdot f(2)=1>0$ αλλά $f(1)=0$.

\vspace{1cm}

%% ===================== Θεώρημα ΜΕΤ =====================
\protash{Θεώρημα Μέγιστης -- Ελάχιστης Τιμής}

Αν μια συνάρτηση $f$ είναι συνεχής σε ένα \textbf{κλειστό} διάστημα $[a,\beta]$, τότε η $f$ παίρνει μέγιστη και ελάχιστη τιμή στο $[a,\beta]$.

\begin{center}
\begin{mytblr}{}
\circled{1}: $A\Rightarrow B$ & \SetCell{orange!90!black}\circled{2}: $B\Rightarrow A$ & \circled{3}: $\neg B\Rightarrow \neg A$ & \SetCell{orange!90!black}\circled{4}: $\neg A\Rightarrow\neg B$\\
\true & \SetCell{orange!50!white}\false & \true & \SetCell{orange!50!white}\false
\end{mytblr}
\end{center}

\theor{\true}
Αν η $f$ είναι \textbf{συνεχής} στο \textbf{κλειστό} διάστημα $[a,\beta]$ $\Rightarrow$ Η $f$ παίρνει μέγιστη και ελάχιστη τιμή στο $[a,\beta]$.\\\\
\antistr{\false}
{\kerkissans 1η μορφή:} Αν η $f$ παίρνει μέγιστη και ελάχιστη τιμή στο $[a,\beta]$ $\Rightarrow$ Η $f$ είναι συνεχής στο $[a,\beta]$.\\
\textcolor{error}{\faTimes\ \textbf{Αντιπαράδειγμα:}} $f(x) = \begin{cases} 1, & x=0 \\ 0, & x\ne 0 \end{cases}$ στο $[-1,1]$: παίρνει $\max{}=1$, $\min{}=0$ αλλά ασυνεχής στο $0$.\\\\
{\kerkissans 2η μορφή:} Αν η $f$ είναι συνεχής σε ένα \textbf{ανοιχτό} διάστημα $(a,\beta)$ $\Rightarrow$ Η $f$ παίρνει μέγιστη και ελάχιστη τιμή.\\
\textcolor{error}{\faTimes\ \textbf{Αντιπαράδειγμα:}} $f(x) = x$ στο $(0,1)$: συνεχής αλλά δεν παίρνει ούτε μέγιστη ούτε ελάχιστη τιμή (πλησιάζει αλλά δεν φτάνει τα $0$ και $1$).\\\\
\aath{\true}
Αν η $f$ είναι \textbf{συνεχής} στο $[a,\beta]$ και \textbf{δεν} παίρνει μέγιστη ή ελάχιστη τιμή $\Rightarrow$ Αδύνατο.\\
{\small Ισοδύναμο με το θεώρημα.}\\\\
\aaa{\false}
Αν η $f$ \textbf{δεν} είναι συνεχής στο $[a,\beta]$ $\Rightarrow$ Η $f$ \textbf{δεν} παίρνει μέγιστη ή ελάχιστη τιμή.\\
\textcolor{error}{\faTimes\ \textbf{Αντιπαράδειγμα:}} Πάλι η $f(x) = \begin{cases} 1, & x=0 \\ 0, & x\ne 0 \end{cases}$: ασυνεχής αλλά έχει μέγιστο και ελάχιστο.

\vspace{1cm}

%% ===================== Παράγωγος & Συνέχεια =====================
\protash{Παράγωγος και συνέχεια}

Αν μια συνάρτηση $f$ είναι παραγωγίσιμη σε ένα σημείο $x_0$, τότε η $f$ είναι και συνεχής στο $x_0$.

\begin{center}
\begin{mytblr}{}
\circled{1}: $A\Rightarrow B$ & \SetCell{orange!90!black}\circled{2}: $B\Rightarrow A$ & \circled{3}: $\neg B\Rightarrow \neg A$ & \SetCell{orange!90!black}\circled{4}: $\neg A\Rightarrow\neg B$\\
\true & \SetCell{orange!50!white}\false & \true & \SetCell{orange!50!white}\false
\end{mytblr}
\end{center}

\theor{\true}
Αν η $f$ είναι παραγωγίσιμη στο $x_0$ $\Rightarrow$ Η $f$ είναι συνεχής στο $x_0$.\\\\
\antistr{\false}
Αν η $f$ είναι συνεχής στο $x_0$ $\Rightarrow$ Η $f$ είναι παραγωγίσιμη στο $x_0$.\\
\textcolor{error}{\faTimes\ \textbf{Αντιπαράδειγμα:}} $f(x)=|x|$: συνεχής στο $0$ αλλά η $f'(0)$ δεν ορίζεται.

\aath{\true}
Αν η $f$ \textbf{δεν} είναι συνεχής στο $x_0$ $\Rightarrow$ Η $f$ \textbf{δεν} είναι παραγωγίσιμη στο $x_0$.\\\\
\aaa{\false}
Αν η $f$ \textbf{δεν} είναι παραγωγίσιμη στο $x_0$ $\Rightarrow$ Η $f$ \textbf{δεν} είναι συνεχής στο $x_0$.\\
\textcolor{error}{\faTimes\ \textbf{Αντιπαράδειγμα:}} Πάλι η $f(x)=|x|$: δεν παραγωγίσιμη στο $0$ αλλά \textbf{είναι} συνεχής.

\vspace{1cm}

%% ===================== Rolle =====================
\protash{Θεώρημα \eng{Rolle}}

Αν μια συνάρτηση $f$ είναι 
\begin{multicols}{3}
\begin{itemize}
\item συνεχής στο $[a,\beta]$, 
\item παραγωγίσιμη στο $(a,\beta)$ και
\item $f(a)=f(\beta)$
\end{itemize}
\end{multicols}
τότε υπάρχει τουλάχιστον ένα $\xi\in(a,\beta)$ τέτοιο ώστε $f'(\xi)=0$.

\begin{center}
\begin{mytblr}{}
\circled{1}: $A\Rightarrow B$ & \SetCell{orange!90!black}\circled{2}: $B\Rightarrow A$ & \circled{3}: $\neg B\Rightarrow \neg A$ & \SetCell{orange!90!black}\circled{4}: $\neg A\Rightarrow\neg B$\\
\true & \SetCell{orange!50!white}\false & \true & \SetCell{orange!50!white}\false
\end{mytblr}
\end{center}

\theor{\true}
Αν η $f$ είναι \textbf{(α)} συνεχής στο $[a,\beta]$, \textbf{(β)} παραγωγίσιμη στο $(a,\beta)$ και \textbf{(γ)} ισχύει $f(a)=f(\beta)$\\
$\Rightarrow$ Υπάρχει τουλάχιστον ένα $\xi\in(a,\beta)$ τέτοιο ώστε $f'(\xi)=0$.\\
{\small Γεωμετρικά: υπάρχει σημείο με \textbf{οριζόντια εφαπτομένη}.}\\\\
\antistr{\false}
{\kerkissans 1η μορφή:} Αν υπάρχει $\xi\in(a,\beta)$ τέτοιο ώστε $f'(\xi)=0$ $\Rightarrow$ Ισχύει $f(a)=f(\beta)$.\\
\textcolor{error}{\faTimes\ \textbf{Αντιπαράδειγμα:}} $f(x)=x^2$ στο $[-1,2]$: $f'(0)=0$ αλλά $f(-1)=1\ne 4=f(2)$.\\\\
{\kerkissans 2η μορφή:} Αν η $f$ είναι συνεχής και ισχύει $f(a)=f(\beta)$ $\Rightarrow$ Υπάρχει $\xi\in(a,\beta)$ τέτοιο ώστε $f'(\xi)=0$.\\
\textcolor{error}{\faTimes\ \textbf{Πρόβλημα:}} Λείπει η «παραγωγίσιμη». Π.χ.\ $f(x)=|x|$ στο $[-1,1]$.\\\\
\aath{\true}
Αν η $f$ είναι συνεχής, παραγωγίσιμη και ισχύει $f'(\xi)\ne 0$ για κάθε $\xi\in(a,\beta)$ $\Rightarrow$ Ισχύει $f(a)\ne f(\beta)$.\\
{\small Αν δεν μηδενίζεται η $f'$, τα άκρα δεν ισούνται.}\\\\
\aaa{\false}
Αν η $f$ είναι συνεχής, παραγωγίσιμη και ισχύει $f(a)\ne f(\beta)$ $\Rightarrow$ Ισχύει $f'(\xi)\ne 0$ για κάθε $\xi\in(a,\beta)$.\\
\textcolor{error}{\faTimes\ \textbf{Αντιπαράδειγμα:}} $f(x) = x^3 - 3x$ στο $[0,3]$: $f(0)\ne f(3)$ αλλά $f'(1)=0$.

\vspace{1cm}

%% ===================== ΘΜΤ =====================
\protash{Θεώρημα Μέσης Τιμής (ΘΜΤ)}

Αν μια συνάρτηση $f$ είναι 
\begin{multicols}{2}
\begin{itemize}
\item συνεχής στο $[a,\beta]$ και 
\item παραγωγίσιμη στο $(a,\beta)$,
\end{itemize}
\end{multicols}
τότε υπάρχει τουλάχιστον ένα $\xi\in(a,\beta)$ τέτοιο ώστε $f'(\xi)=\dfrac{f(\beta)-f(a)}{\beta-a}$.

\begin{center}
\begin{mytblr}{}
\circled{1}: $A\Rightarrow B$ & \SetCell{orange!90!black}\circled{2}: $B\Rightarrow A$ & \circled{3}: $\neg B\Rightarrow \neg A$ & \SetCell{orange!90!black}\circled{4}: $\neg A\Rightarrow\neg B$\\
\true & \SetCell{orange!50!white}\false & \true & \SetCell{orange!50!white}\false
\end{mytblr}
\end{center}

\theor{\true}
Αν η $f$ είναι \textbf{(α)} συνεχής στο $[a,\beta]$ και \textbf{(β)} παραγωγίσιμη στο $(a,\beta)$\\
$\Rightarrow$ Υπάρχει τουλάχιστον ένα $\xi\in(a,\beta)$ τέτοιο ώστε $f'(\xi)=\dfrac{f(\beta)-f(a)}{\beta - a}$.\\
{\small Γεωμετρικά: η εφαπτομένη στο $M(\xi,f(\xi))$ είναι παράλληλη στη χορδή $AB$.}\\\\
\antistr{\false}
Αν υπάρχει $\xi\in(a,\beta)$ τέτοιο ώστε $f'(\xi) = \frac{f(\beta)-f(a)}{\beta-a}$ $\Rightarrow$ Η $f$ είναι συνεχής στο $[a,\beta]$.\\
\textcolor{error}{\faTimes\ \textbf{Πρόβλημα:}} Παραγωγίσιμη στο ανοιχτό δεν εγγυάται συνέχεια στα \textbf{κλειστά} άκρα.\\\\
\aath{\true}
Αν η $f$ είναι συνεχής, παραγωγίσιμη και \textbf{δεν} υπάρχει $\xi\in(a,\beta)$ τέτοιο ώστε $f'(\xi)=\frac{f(\beta)-f(a)}{\beta-a}$ $\Rightarrow$ Αδύνατο.\\
{\small Η άρνηση του συμπεράσματος + υποθέσεις = άτοπο.}\\\\
\aaa{\false}
Αν η $f$ είναι συνεχής και παραγωγίσιμη $\Rightarrow$ Υπάρχει \textbf{μοναδικό} $\xi\in(a,\beta)$ τέτοιο ώστε $f'(\xi)=\frac{f(\beta)-f(a)}{\beta-a}$.\\
\textcolor{error}{\faTimes\ \textbf{Αντιπαράδειγμα:}} $f(x)=\hm{x}$ στο $[0,2\pi]$: $f'(\xi)=\syn{\xi}=0$ σε \textbf{δύο} σημεία.

\vspace{1cm}

%% ===================== Fermat =====================
\protash{Θεώρημα \eng{Fermat} (τοπικό ακρότατο)}

Αν μια συνάρτηση $f$ είναι παραγωγίσιμη σε ένα σημείο $x_0$ και παρουσιάζει τοπικό ακρότατο σε αυτό, τότε $f'(x_0)=0$.

\begin{center}
\begin{mytblr}{}
\circled{1}: $A\Rightarrow B$ & \SetCell{orange!90!black}\circled{2}: $B\Rightarrow A$ & \circled{3}: $\neg B\Rightarrow \neg A$ & \SetCell{orange!90!black}\circled{4}: $\neg A\Rightarrow\neg B$\\
\true & \SetCell{orange!50!white}\false & \true & \SetCell{orange!50!white}\false
\end{mytblr}
\end{center}

\theor{\true}
Αν η $f$ είναι παραγωγίσιμη στο $x_0$ και παρουσιάζει τοπικό ακρότατο στο $x_0$ $\Rightarrow$ Ισχύει $f'(x_0)=0$.\\
{\small Γεωμετρικά: σε κορυφή/κοιλάδα η εφαπτομένη είναι \textbf{οριζόντια}.}\\\\
\antistr{\false}
{\kerkissans 1η μορφή:} Αν ισχύει $f'(x_0) = 0$ $\Rightarrow$ Η $f$ παρουσιάζει τοπικό ακρότατο στο $x_0$.\\
\textcolor{error}{\faTimes\ \textbf{Κλασική παγίδα.}} $f(x)=x^3$: $f'(0)=0$ αλλά, δεν αποτελεί ακρότατο ($f\nearrow$).
{\kerkissans 2η μορφή:} Αν η $f$ παρουσιάζει τοπικό ακρότατο στο $x_0$ $\Rightarrow$ Ισχύει $f'(x_0) = 0$.\\
\textcolor{error}{\faTimes\ \textbf{Πρόβλημα:}} Λείπει η συνθήκη \textbf{παραγωγίσιμη}. Π.χ.Η $f(x)=|x|$ έχει ακρότατο αλλά η $f'(0)$ δεν ορίζεται.\\\\
\aath{\true}
Αν η $f$ είναι παραγωγίσιμη στο $x_0$ και ισχύει $f'(x_0)\ne 0$ $\Rightarrow$ Η $f$ \textbf{δεν} παρουσιάζει τοπικό ακρότατο στο $x_0$.\\\\
\aaa{\false}
Αν η $f$ \textbf{δεν} παρουσιάζει τοπικό ακρότατο στο $x_0$ $\Rightarrow$ Ισχύει $f'(x_0)\ne 0$.\\
\textcolor{error}{\faTimes\ \textbf{Αντιπαράδειγμα:}} $f(x) = x^3$: δεν έχει ακρότατο στο $0$ αλλά $f'(0) = 0$.

\vspace{1cm}

%% ===================== Κριτήριο γνησίως αύξουσας =====================
\protash{Κριτήριο γνησίως αύξουσας (πρόσημο $f'$)}

Αν μια συνάρτηση $f$ είναι συνεχής σε ένα διάστημα $\Delta$ και ισχύει $f'(x)>0$ σε κάθε εσωτερικό σημείο του, τότε η $f$ είναι γνησίως αύξουσα στο $\Delta$.

\begin{center}
\begin{mytblr}{}
\circled{1}: $A\Rightarrow B$ & \SetCell{orange!90!black}\circled{2}: $B\Rightarrow A$ & \circled{3}: $\neg B\Rightarrow \neg A$ & \SetCell{orange!90!black}\circled{4}: $\neg A\Rightarrow\neg B$\\
\true & \SetCell{orange!50!white}\false & \true & \SetCell{orange!50!white}\false
\end{mytblr}
\end{center}

\theor{\true}
Αν η $f$ είναι συνεχής στο $\Delta$ και ισχύει $f'(x)>0$ σε κάθε εσωτερικό σημείο του $\Rightarrow$ Η $f$ είναι γνησίως αύξουσα στο $\Delta$.\\
{\small Σημείωση: Ισχύει και με $f'(x)\ge 0$ αν οι ρίζες είναι μεμονωμένες.}\\\\
\antistr{\false}
Αν η $f$ είναι γνησίως αύξουσα $\Rightarrow$ Ισχύει $f'(x)>0$ σε κάθε εσωτερικό σημείο του $\Delta$.\\
\textcolor{error}{\faTimes\ \textbf{Κλασική παγίδα.}} $f(x)=x^3$: γν. αύξουσα αλλά $f'(0)=0$. Το σωστό είναι: $f'(x)\ge 0$.\\\\
\aath{\true}
Αν η $f$ \textbf{δεν} είναι γνησίως αύξουσα $\Rightarrow$ Η ανισότητα $f'(x)>0$ \textbf{δεν} ισχύει σε κάθε εσωτερικό σημείο του $\Delta$.\\\\
\aaa{\false}
Αν η ανισότητα $f'(x)>0$ \textbf{δεν} ισχύει σε κάθε εσωτερικό σημείο $\Rightarrow$ Η $f$ \textbf{δεν} είναι γνησίως αύξουσα.\\
\textcolor{error}{\faTimes\ \textbf{Αντιπαράδειγμα:}} $f(x) = x^3$: $f'(0)=0$ αλλά αύξουσα.

\vspace{1cm}

%% ===================== Κριτήριο γνησίως φθίνουσας =====================
\protash{Κριτήριο γνησίως φθίνουσας (πρόσημο $f'$)}

Αν μια συνάρτηση $f$ είναι συνεχής σε ένα διάστημα $\Delta$ και ισχύει $f'(x)<0$ σε κάθε εσωτερικό σημείο του, τότε η $f$ είναι γνησίως φθίνουσα στο $\Delta$.

\begin{center}
\begin{mytblr}{}
\circled{1}: $A\Rightarrow B$ & \SetCell{orange!90!black}\circled{2}: $B\Rightarrow A$ & \circled{3}: $\neg B\Rightarrow \neg A$ & \SetCell{orange!90!black}\circled{4}: $\neg A\Rightarrow\neg B$\\
\true & \SetCell{orange!50!white}\false & \true & \SetCell{orange!50!white}\false
\end{mytblr}
\end{center}

\theor{\true}
Αν η $f$ είναι συνεχής στο $\Delta$ και ισχύει $f'(x)<0$ σε κάθε εσωτερικό σημείο του $\Rightarrow$ Η $f$ είναι γνησίως φθίνουσα στο $\Delta$.\\\\
\antistr{\false}
Αν η $f$ είναι γνησίως φθίνουσα $\Rightarrow$ Ισχύει $f'(x)<0$ σε κάθε εσωτερικό σημείο του $\Delta$.\\
\textcolor{error}{\faTimes\ \textbf{Αντιπαράδειγμα:}} $f(x)=-x^3$: γν. φθίνουσα αλλά $f'(0)=0$. Το σωστό είναι: $f'(x)\le 0$.\\\\
\aath{\true}
Αν η $f$ \textbf{δεν} είναι γνησίως φθίνουσα $\Rightarrow$ Η ανισότητα $f'(x)<0$ \textbf{δεν} ισχύει σε κάθε εσωτερικό σημείο του $\Delta$.\\\\
\aaa{\false}
Αν η ανισότητα $f'(x)<0$ \textbf{δεν} ισχύει σε κάθε εσωτερικό σημείο $\Rightarrow$ Η $f$ \textbf{δεν} είναι γνησίως φθίνουσα.\\
\textcolor{error}{\faTimes\ \textbf{Αντιπαράδειγμα:}} $f(x) = -x^3$: $f'(0)=0$ αλλά φθίνουσα.

\vspace{1cm}

%% ===================== Θεώρημα Σταθερής Συνάρτησης =====================
\protash{Θεώρημα Σταθερής Συνάρτησης}

Αν μια συνάρτηση $f$ είναι συνεχής σε ένα διάστημα $\Delta$ και ισχύει $f'(x)=0$ σε κάθε εσωτερικό σημείο του, τότε η $f$ είναι σταθερή στο $\Delta$.

\begin{center}
\begin{mytblr}{}
\circled{1}: $A\Rightarrow B$ & \SetCell{orange!90!black}\circled{2}: $B\Rightarrow A$ & \circled{3}: $\neg B\Rightarrow \neg A$ & \SetCell{orange!90!black}\circled{4}: $\neg A\Rightarrow\neg B$\\
\true & \SetCell{orange!50!white}\true & \true & \SetCell{orange!50!white}\true
\end{mytblr}
\end{center}

{\small \textbf{Ιδιαίτερο:} Εδώ όλα ισχύουν (ισοδυναμία). $f' = 0$ σε \textbf{διάστημα} $\Leftrightarrow$ $f$ σταθερή.}\\\\
\theor{\true}
Αν η $f$ είναι συνεχής στο \textbf{διάστημα} $\Delta$ και ισχύει $f'(x)=0$ σε κάθε εσωτερικό σημείο του $\Rightarrow$ Η $f$ είναι σταθερή στο $\Delta$.\\\\
\antistr{\true}
Αν η $f$ είναι σταθερή στο $\Delta$ $\Rightarrow$ Ισχύει $f'(x)=0$ σε κάθε εσωτερικό σημείο του $\Delta$.\\\\
\aath{\true}
Αν η $f$ \textbf{δεν} είναι σταθερή στο $\Delta$ $\Rightarrow$ Η παράγωγος $f'(x)$ \textbf{δεν} είναι παντού μηδέν στο εσωτερικό του $\Delta$.\\\\
\aaa{\true}
Αν η παράγωγος $f'(x)$ \textbf{δεν} είναι παντού μηδέν στο εσωτερικό του $\Delta$ $\Rightarrow$ Η $f$ \textbf{δεν} είναι σταθερή στο $\Delta$.\\\\
\kerkissans{{\large \textcolor{error}{\faExclamationTriangle}\ Η μεγάλη παγίδα: Όχι διάστημα.}}

Αν $f'(x)=0$ σε κάθε σημείο ενός \textbf{συνόλου} $A$ (όχι διάστημα) $\Rightarrow$ η $f$ \textbf{δεν} είναι αναγκαστικά σταθερή.\\
\textcolor{error}{\faTimes\ \textbf{Αντιπαράδειγμα:}} $f(x)=\frac{x}{|x|}$ στο $\mathbb{R}^*$: $f'(x)=0$ παντού αλλά $f=1$ στο $(0,+\infty)$ και $f=-1$ στο $(-\infty,0)$.

\vspace{1cm}

%% ===================== ΘΕΤ =====================
\protash{Θεώρημα Ενδιάμεσων Τιμών (ΘΕΤ)}

Αν μια συνάρτηση $f$ είναι συνεχής στο $[\alpha,\beta]$ και ισχύει $f(\alpha)\ne f(\beta)$, τότε για κάθε αριθμό $\eta$ μεταξύ $f(\alpha)$ και $f(\beta)$ υπάρχει τουλάχιστον ένα $x_0\in(\alpha,\beta)$ τέτοιο ώστε $f(x_0)=\eta$.

\begin{center}
\begin{mytblr}{}
\circled{1}: $A\Rightarrow B$ & \SetCell{orange!90!black}\circled{2}: $B\Rightarrow A$ & \circled{3}: $\neg B\Rightarrow \neg A$ & \SetCell{orange!90!black}\circled{4}: $\neg A\Rightarrow\neg B$\\
\true & \SetCell{orange!50!white}\false & \true & \SetCell{orange!50!white}\false
\end{mytblr}
\end{center}

\theor{\true}
Αν η $f$ είναι συνεχής στο $[\alpha,\beta]$ και ισχύει $f(\alpha)\ne f(\beta)$
$\Rightarrow$ Για κάθε αριθμό $\eta$ μεταξύ των $f(\alpha)$ και $f(\beta)$, υπάρχει τουλάχιστον ένα $x_0\in(\alpha,\beta)$ τέτοιο ώστε $f(x_0)=\eta$.\\\\
\antistr{\false}
Αν η $f$ παίρνει όλες τις ενδιάμεσες τιμές μεταξύ των $f(\alpha)$ και $f(\beta)$ $\Rightarrow$ Η $f$ είναι συνεχής.\\
\textcolor{error}{\faTimes\ \textbf{Πρόβλημα:}} Υπάρχουν ασυνεχείς συναρτήσεις με αυτή την ιδιότητα.\\\\
\aath{\true}
Αν η $f$ \textbf{δεν} παίρνει κάποια ενδιάμεση τιμή $\eta$ $\Rightarrow$ Η $f$ \textbf{δεν} είναι συνεχής στο $[\alpha,\beta]$.\\\\
\aaa{\false}
Αν η $f$ \textbf{δεν} είναι συνεχής $\Rightarrow$ Η $f$ \textbf{δεν} παίρνει όλες τις ενδιάμεσες τιμές.\\
\textcolor{error}{\faTimes\ \textbf{Αντιπαράδειγμα:}} Οι συναρτήσεις που αναφέραμε στο \circled{2}.

\vspace{1cm}

%% ===================== Κυρτότητα =====================
\protash{Κυρτότητα και $f''$}

Αν μια συνάρτηση $f$ είναι συνεχής σε ένα διάστημα $\Delta$ και ισχύει $f''(x)>0$ στο εσωτερικό του, τότε η $f$ είναι κυρτή (στρέφει τα κοίλα άνω) στο $\Delta$.

\begin{center}
\begin{mytblr}{}
\circled{1}: $A\Rightarrow B$ & \SetCell{orange!90!black}\circled{2}: $B\Rightarrow A$ & \circled{3}: $\neg B\Rightarrow \neg A$ & \SetCell{orange!90!black}\circled{4}: $\neg A\Rightarrow\neg B$\\
\true & \SetCell{orange!50!white}\false & \true & \SetCell{orange!50!white}\false
\end{mytblr}
\end{center}

\theor{\true}
Αν η $f$ είναι συνεχής στο $\Delta$ και ισχύει $f''(x)>0$ στο εσωτερικό του $\Rightarrow$ Η $f$ είναι \textbf{κυρτή} στο $\Delta$.\\\\
\antistr{\false}
Αν η $f$ είναι κυρτή $\Rightarrow$ Ισχύει $f''(x)>0$ παντού στο εσωτερικό του $\Delta$.\\
\textcolor{error}{\faTimes\ \textbf{Αντιπαράδειγμα:}} $f(x)=x^4$: κυρτή παντού αλλά $f''(0)=0$. Το σωστό είναι: $f''(x)\ge 0$.

\aath{\true}
Αν η $f$ \textbf{δεν} είναι κυρτή $\Rightarrow$ Η ανισότητα $f''(x)>0$ \textbf{δεν} μπορεί να ισχύει παντού.\\\\
\aaa{\false}
Αν η ανισότητα $f''(x)>0$ \textbf{δεν} ισχύει παντού $\Rightarrow$ Η $f$ \textbf{δεν} είναι κυρτή.\\
\textcolor{error}{\faTimes\ \textbf{Αντιπαράδειγμα:}} $f(x)=x^4$: $f''(0)=0$ αλλά κυρτή.

\begin{center}
{\LARGE \kerkissans{\textbf{Συγκεντρωτικός πίνακας}}}
\end{center}

\begin{center}
{\small
\begin{mytblr}{}
\textbf{Θεώρημα} & \circled{1} & \circled{2} & \circled{3} & \circled{4}\\
Μονοτονία $\Rightarrow$ $1$-$1$ & \true & \false & \true & \false\\
\eng{Bolzano} & \true & \false & \true & \false\\
Μεγ.--Ελάχ. Τιμής & \true & \false & \true & \false\\
Παραγωγίσιμη $\Rightarrow$ Συνεχής & \true & \false & \true & \false\\
\eng{Rolle} & \true & \false & \true & \false\\
Θ.Μ.Τ. & \true & \false & \true & \false\\
\eng{Fermat} (τοπ. ακρότατο) & \true & \false & \true & \false\\
Κριτήριο γν. αύξουσας & \true & \false & \true & \false\\
Κριτήριο γν. φθίνουσας & \true & \false & \true & \false\\
Σταθερή συνάρτηση & \true & \true & \true & \true\\
Θ.Ε.Τ. & \true & \false & \true & \false\\
Κυρτότητα ($f''>0$) & \true & \false & \true & \false
\end{mytblr}
}
\end{center}
{\small \textbf{Υπενθύμιση:} \circled{1}$\Leftrightarrow$\circled{3} πάντα, και \circled{2}$\Leftrightarrow$\circled{4} πάντα. Το μόνο θεώρημα όπου όλα ισχύουν είναι η \textbf{Σταθερή Συνάρτηση} (ισοδυναμία).}
